\unit{Exposé 182 — Géométrie formelle et géométrie algébrique}
\sid{182}
\src{193}
\seminar{(Mai 1959)}
\paperhead{GÉOMÉTRIE FORMELLE ET GÉOMÉTRIE ALGÉBRIQUE}%
  {par Alexander GROTHENDIECK}

\fsection{1}{Schémas}

On sait qu'un espace algébrique affine défini sur un corps $k$ est
essentiellement déterminé par son algèbre affine $A$ (anneau des fonctions
régulières définies sur $k$), les morphismes d'espaces algébriques
$X\longrightarrow Y$ correspondant biunivoquement aux homomorphismes de
$k$-algèbres $A(Y)\longrightarrow A(X)$. L'algèbre affine correspondant à un
espace algébrique est une $k$-algèbre de type fini et, dans le point de vue
« classique », elle n'a pas d'éléments nilpotents ; inversement, toute algèbre
de ce type est obtenue comme algèbre affine d'un espace algébrique défini sur
$k$. Il y a alors un dictionnaire connu permettant d'interpréter les situations
concernant des espaces algébriques affines en termes d'algèbre commutative. On
a constaté depuis longtemps qu'on obtenait alors des énoncés plus généraux, car
il n'était généralement plus besoin de supposer que les anneaux en jeu étaient
du type juste envisagé, l'hypothèse noethérienne étant le plus souvent
suffisante. En particulier, qu'il y ait ou non un corps de base donné, il n'y
avait pas lieu d'exclure le cas où ces anneaux contiennent des éléments
nilpotents. Jusqu'à présent, les géomètres s'étaient refusés à tenir compte de
ces indications et se sont obstinés à se restreindre à la considération
d'algèbres affines sans éléments nilpotents, i.e. d'espaces algébriques dans les
faisceaux structuraux desquels il n'y a pas d'éléments nilpotents (et même le
plus souvent, des espaces algébriques « absolument irréductibles »). Le
conférencier pense que cet état d'esprit a été un obstacle sérieux au
développement des méthodes vraiment naturelles en Géométrie algébrique.

Soit $A$ un anneau commutatif. Il est bien connu que l'ensemble
$X=\operatorname{Spec}(A)$ des idéaux premiers de $A$ est muni d'une topologie
naturelle, la « topologie de Zariski » ou topologie spectrale. D'autre part, il
y a un faisceau d'anneaux commutatifs $\Osh_X$ sur $X$, dont la fibre en
$\mathfrak p\in X$ est l'anneau localisé $A_{\mathfrak p}$, et dont l'anneau des
sections s'identifie à $A$. Ainsi, $X$ devient un \emph{espace annelé}, appelé
\emph{spectre premier} de $A$. Un homomorphisme d'anneaux
$f:A\longrightarrow B$ définit un morphisme d'espaces annelés
$f':\operatorname{Spec}(B)\longrightarrow\operatorname{Spec}(A)$,
l'application ensembliste sous-jacente n'étant autre que
$\mathfrak p\longmapsto f^{-1}(\mathfrak p)$. Les homomorphismes

\src{194}
d'espaces annelés de $\operatorname{Spec}(B)$ dans
$\operatorname{Spec}(A)$ obtenus de cette façon sont exactement ceux pour
lesquels les homomorphismes $\Osh_x\longrightarrow\Osh_y$ ($x=f'(y)$) sont
locaux (i.e. l'image inverse de l'idéal maximal est l'idéal maximal).

On appelle \emph{schéma affine} un espace annelé isomorphe à un
$\operatorname{Spec}(A)$, et \emph{préschéma} un espace annelé localement
affine, i.e. dont tout point a un voisinage ouvert qui est un schéma affine pour
la structure induite. On définit de façon évidente les \emph{morphismes} des
préschémas ; localement ils correspondent à des homomorphismes d'anneaux.

Quand on se fixe un préschéma $S$, et qu'on regarde des morphismes de préschémas
$X\longrightarrow S$, alors $S$ joue le rôle d'un corps ou d'un anneau de base
(ou mieux, d'un espace de base dans une fibration). On dit alors que $X$ est un
$S$-préschéma ; si $S=\operatorname{Spec}(A)$, cela signifie aussi que $\Osh_X$
est un faisceau de $A$-algèbres. Ainsi, tout préschéma peut être regardé de
façon unique comme un $\mathbf Z$-préschéma. Bien entendu, les $S$-préschémas
forment une catégorie ; de plus on montre que dans cette catégorie le produit
de deux objets $X,Y$ existe toujours, il est noté $X\times_S Y$. Cette notion
de produit permet de définir le changement de base dans un $S$-préschéma,
correspondant à un morphisme $S'\longrightarrow S$ : en effet,
$X\times_S S'$ pourra être considéré comme un $S'$-préschéma.

On dit que $X$ est \emph{séparé} au-dessus de $S$ si la diagonale de
$X\times_S X$ est fermée. On appelle \emph{schéma} un préschéma séparé
au-dessus de $\mathbf Z$ ; il est alors séparé au-dessus de n'importe quoi.
Pour simplifier, nous ne parlerons plus que de schémas, que de plus nous
supposerons \emph{noethériens}, i.e. réunions finies d'ouverts affines, spectres
d'anneaux noethériens. $X$ est dit \emph{de type fini sur} $S$ si, pour tout
ouvert affine $U$ de $S$, son image inverse dans $X$ est réunion finie d'ouverts
affines dont les anneaux sont des algèbres de type fini sur l'anneau de $U$. Ce
sont de tels $S$-schémas qui se prêtent à une étude proprement géométrique. En
particulier, pour tout $s\in S$, la fibre $f^{-1}(s)$ de $X$ au-dessus de $s$
est un schéma algébrique sur le corps résiduel $\kappa(s)$ de l'anneau local
$\Osh_s$ de $s$ dans $S$. Ainsi, $X$ peut dans une certaine mesure être
considéré comme une famille d'« espaces algébriques » $f^{-1}(s)$, le paramètre
$s$ parcourant $S$ (i.e., du point de vue local, l'ensemble des idéaux premiers
d'un anneau donné). Bien entendu, les $\kappa(s)$ peuvent avoir des
caractéristiques différentes. Si $S=\operatorname{Spec}(k)$, où $k$ est un
corps, on retrouve essentiellement la notion usuelle d'« espace algébrique »,
avec la seule différence que maintenant le faisceau structural peut avoir des
éléments nilpotents.

\src{195}
En s'inspirant de notions bien connues, on définit la notion de
\emph{morphisme projectif}, et plus généralement de \emph{morphisme propre}. Un
tel morphisme est de type fini ; de plus il transforme parties fermées en
parties fermées, et garde cette propriété par changement de base quelconque.

$X$ étant un schéma (noethérien, comme toujours), le faisceau $\Osh_X$ est un
\emph{faisceau cohérent d'anneaux} au sens de [2]. Les faisceaux cohérents de
modules sur $X$ sont donc aussi les faisceaux qui localement sont isomorphes à
un conoyau d'un homomorphisme $\Osh_X^{m}\longrightarrow\Osh_X^{n}$.

\fsection{2}{Schémas formels}

Soient $X$ un schéma, et $X'$ une partie fermée de $X$. Alors il existe un
sous-faisceau cohérent $J$ de $\Osh_X$ tel que
$X'=\operatorname{Supp}(\Osh_X/J)$ (et il en existe même un plus grand). Muni
de $\Osh_X/J$, $X'$ devient un schéma, noté $X_0$ ; un tel schéma est appelé
« sous-schéma fermé de $X$ ». On peut aussi pour tout $n$ considérer $X'$ muni
de $\Osh_X/J^{n+1}$, noté $X_n$ ; c'est un sous-préschéma fermé de $X$ dont
l'ensemble sous-jacent est encore $X'$, mais ayant un autre faisceau
structural, soit $\Osh_{X_n}=\Osh_X/J^{n+1}$. Évidemment les $\Osh_{X_n}$
forment un système projectif de faisceaux d'anneaux sur $X'$, dont la limite
projective $\overline{\Osh}_X$ est appelée \emph{complété formel de $\Osh_X$ le
long de $X'$}. Muni de ce faisceau d'anneaux, $X'$ est appelé
\emph{complété formel de $X$ le long de $X'$}, c'est donc un espace annelé,
mais pas un schéma en général. Pour tout faisceau cohérent $F$ sur $X$, on peut
considérer de même le complété formel
$\overline F=\varprojlim_n F_n$ de $F$ le long de $X'$ (où
$F_n=F\otimes_{\Osh_X}\Osh_X/J^{n+1}$) ; c'est un faisceau de modules sur
$\overline X$. Ses sections s'appellent les « sections formelles de $F$ le long
de $X'$ » et s'identifient aux éléments de
$\varprojlim_n\Gamma(X',F_n)$. Pour $F=\Osh_X$, on trouve les « fonctions
holomorphes » de $X$ le long de $X'$ au sens de Zariski (dont nous ne suivrons
pas la terminologie à cause de ses interférences avec la terminologie
classique).

On appelle \emph{schéma formel} (sous-entendu : noethérien) un espace
topologique $\mathfrak X$, muni d'un faisceau d'anneaux topologiques
$\Osh_{\mathfrak X}$ satisfaisant à la condition suivante : on a un
isomorphisme de faisceaux d'anneaux topologiques
$\Osh_{\mathfrak X}=\varprojlim_n\Osh_n$, où les $\Osh_n$ forment un système
projectif de faisceaux d'anneaux sur $\mathfrak X$, faisant chacun de
$\mathfrak X$ un schéma $\mathfrak X_n$, et tel que pour $m\geq n$,
l'homomorphisme $\Osh_m\longrightarrow\Osh_n$ soit surjectif et ait pour noyau
$J_m^{n+1}$, où $J_m$ est

\src{196}
le noyau de $\Osh_m\longrightarrow\Osh_0$. On montre alors que
$\Osh_{\mathfrak X}$ est un \emph{faisceau cohérent d'anneaux locaux
noethériens}.

En vertu des définitions, un complété formel $\overline X$ comme plus haut est
un schéma formel, et inversement, tout schéma formel est \emph{localement} de ce
type. En fait, la donnée d'un schéma formel \emph{affine} (i.e. tel que
$\mathfrak X_0$ soit affine, ce qui implique que tous les $\mathfrak X_n$ le
sont) est équivalente à la donnée d'un anneau topologique noethérien,
$J$-adique séparé et complet.

Les définitions habituelles : morphisme, morphisme de type fini, morphisme
propre, etc., pour les schémas ordinaires, s'étendent sans difficulté aux
schémas formels.

\fsection{3}{Les trois théorèmes fondamentaux}

Soit $f:X\longrightarrow Y$ un morphisme propre de schémas (noethériens comme
toujours), soit $Y'$ une partie fermée de $Y$, $X'$ son image inverse dans $X$ ;
considérons les complétés formels $\overline Y$, $\overline X$ correspondants.
Alors $f$ induit un morphisme
$\overline f:\overline X\longrightarrow\overline Y$ de schémas formels, qui est
d'ailleurs propre. Soit $F$ un faisceau cohérent sur $X$, alors $\overline F$
est un faisceau cohérent sur $\overline X$. Dans le théorème 1, on oublie
$X,Y,F$ et on ne regarde que le morphisme propre $\overline f$ de schémas
formels, et le faisceau cohérent $\overline F$ sur $\overline X$. (Cependant,
le conférencier n'a écrit de démonstration complète que dans le cas où on part
de $X,Y,f,F$.)

\result{THÉORÈME}{1 (théorème de finitude)}
\begin{enumerate}
  \item[\textit{i.}] Les $R^{q}\overline f_{*}(\overline F)$ sont des faisceaux
    cohérents sur $\overline Y$.
  \item[\textit{ii.}] Les homomorphismes naturels
    \[
      R^{q}\overline f_{*}(\overline F)
      \longrightarrow \varprojlim_n R^{q}f_{n*}(F_n)
    \]
    sont des isomorphismes.
\end{enumerate}

Dans cet énoncé, on suppose choisi un sous-faisceau cohérent $J$ de $\Osh_Y$
définissant $Y'$, d'où par image réciproque un sous-faisceau cohérent de
$\Osh_X$ définissant $X'$, d'où par suite la définition des $F_n,X_n,Y_n$ et
$f_n:X_n\longrightarrow Y_n$ comme au n\textsuperscript{o} 2. Les changements
mineurs à faire dans l'explicitation de ces notations sont évidents, si on
partait d'un morphisme propre quelconque de deux schémas formels.

Le théorème 1 ne concernait que la « cohomologie formelle ». Le théorème suivant

\src{197}
la met en rapport avec la « cohomologie algébrique », et s'apparente à un
théorème bien connu de Serre [4] sur la comparaison entre cohomologie algébrique
et cohomologie analytique.

\result{THÉORÈME}{2 (Premier théorème de comparaison)} Les
$R^{q}f_{*}(F)$ sont des faisceaux cohérents sur $Y$ (ce qui est un cas
particulier du théorème 1), et les homomorphismes naturels
\[
  \overline{R^{q}f_{*}(F)}
  \longrightarrow \varprojlim_n R^{q}f_{n*}(F_n)
\]
sont des isomorphismes.

\result{COROLLAIRE}{1} On a des isomorphismes canoniques :
\[
  \overline{R^{q}f_{*}(F)}=R^{q}\overline f_{*}(\overline F).
\]

Ce corollaire est, pour $q=0$, une généralisation du « théorème fondamental des
fonctions holomorphes » de Zariski, dont nous déduirons une généralisation du
« théorème de connexion » de Zariski. Notons d'ailleurs que, alors que le
théorème 1 (\textit{ii}) est trivial pour $q=0$, il n'en est plus du tout de
même pour le théorème 2 ou pour sa formulation équivalente (corollaire 1). En
fait, la démonstration procède par récurrence descendante sur $q$ (étant
triviale pour $q$ grand, car alors les deux membres sont nuls), et le cas $q=0$
apparaît donc comme le dernier pas de la récurrence, donc si on peut dire le cas
« le plus difficile ».

\result{COROLLAIRE}{2} Supposons $Y=\operatorname{Spec}(A)$, $Y'$ étant défini
par l'idéal $J$ de $A$. Alors, pour tout faisceau cohérent $F$ sur $X$, les
$H^{q}(X,F)$ sont des $A$-modules de type fini, dont les complétés $J$-adiques
sont les $H^{q}(\overline X,\overline F)$.

Appliquant enfin ce corollaire à
$H=\Homsh_{\Osh_X}(F,G)$, on trouve :

\result{COROLLAIRE}{3} Supposons $Y=\operatorname{Spec}(A)$, $Y'$ étant défini
par l'idéal $J$ de $A$. Soient $F,G$ deux faisceaux cohérents sur $X$, alors
$\operatorname{Hom}(F,G)$ est un module de type fini sur $A$, dont le complété
$J$-adique s'identifie à
$\operatorname{Hom}(\overline F,\overline G)$.

Bien entendu, l'application naturelle
$\operatorname{Hom}(F,G)\longrightarrow
\operatorname{Hom}(\overline F,\overline G)$ est celle qui associe à un
homomorphisme $u:F\longrightarrow G$ son prolongement « par continuité »
$\overline u:\overline F\longrightarrow\overline G$ (moyennant quoi
$\overline F$ devient un foncteur en $F$).

Supposons maintenant que $A$ soit séparé et complet pour sa topologie
$J$-adique, alors les corollaires 2 et 3 précédents donnent :
\[
  H^{q}(X,F)=H^{q}(\overline X,\overline F),
  \qquad
  \operatorname{Hom}(F,G)
  =\operatorname{Hom}(\overline F,\overline G).
\]

\src{198}
Cette dernière identité montre que la catégorie des faisceaux cohérents sur
$X$ s'identifie à une \emph{sous-catégorie} (avec comme morphismes les
morphismes induits) de la catégorie des faisceaux cohérents sur $\overline X$.
En fait, on a même :

\result{THÉORÈME}{3} Pour qu'un faisceau de modules sur $\overline X$ soit
cohérent, il faut et il suffit qu'il soit isomorphe à un faisceau de la forme
$\overline F$, où $F$ est un faisceau cohérent sur $X$ (déterminé à un
isomorphisme canonique près en vertu du théorème 2, corollaire 3). [On rappelle
que maintenant $Y=\operatorname{Spec}(A)$, $A$ étant un anneau topologique
noethérien $J$-adique séparé et complet.]

\result{COROLLAIRE}{1} Les sous-schémas fermés de $X$ correspondent
biunivoquement aux sous-schémas formels fermés de $\overline X$.

En effet, ils correspondent aux sous-faisceaux cohérents de $\Osh_X$, resp. de
$\overline{\Osh}_X$. Regardant les graphes de morphismes comme des sous-schémas
fermés, on déduit du corollaire 1 :

\result{COROLLAIRE}{2} Soient $X,Z$ deux schémas propres au-dessus de $A$
(anneau noethérien $J$-adique, séparé et complet). Alors l'application
$g\longmapsto\overline g$ définit une correspondance biunivoque entre les
$Y$-morphismes de $X$ dans $Z$, et les $\overline Y$-morphismes de
$\overline X$ dans $\overline Z$.

En d'autres termes, les schémas algébriques propres sur $A$ apparaissent comme
une \emph{sous-catégorie} (avec comme morphismes les morphismes induits) de la
catégorie des schémas formels propres sur $\overline Y$. On fera attention
cependant qu'il existe des schémas formels propres sur $\overline Y$ qui ne
sont pas « algébrisables », i.e. isomorphes à un $\overline X$, avec $X$ propre
sur $A$ (tout comme il existe des variétés analytiques complexes compactes qui
ne proviennent pas de variétés algébriques définies sur le corps des complexes).
De tels schémas formels s'introduisent de façon naturelle en « théorie des
modules ». Notons cependant un cas particulier intéressant où un schéma formel
est algébrisable :

\result{THÉORÈME}{4} Soit $A$ un anneau noethérien local complet de corps
résiduel $k$, et soit $\mathfrak X$ un schéma formel propre sur $A$ (muni de sa
topologie $\mathfrak m(A)$-adique). On suppose :
\begin{enumerate}
  \item[\textit{i.}] que les anneaux locaux de $\Osh_{\mathfrak X}$ sont des
    modules plats sur $A$, ou ce qui revient au même, si on munit
    $\Osh_{\mathfrak X}$ et $A$ de la filtration définie par les puissances de
    l'idéal maximal de $A$, on a pour les gradués associés la relation
    \[
      \operatorname{gr}(\Osh_{\mathfrak X})
      \simeq
      \operatorname{gr}^{0}(\Osh_{\mathfrak X})
      \otimes_k\operatorname{gr}(A).
    \]
  \item[\textit{ii.}] Considérant
    $\mathfrak X_0=\mathfrak X\otimes_A k$ comme un schéma algébrique sur $k$,
    on a
    \[
      H^2(\mathfrak X_0,\Osh_{\mathfrak X_0})=0.
    \]
  \item[\textit{iii.}] $\mathfrak X_0$ est projectif.
\end{enumerate}
Sous ces conditions, $\mathfrak X$ est algébrisable, et de façon précise est
isomorphe à $\overline X$, où $X$ est un $A$-schéma projectif.

Les conditions \textit{(ii)} et \textit{(iii)} seront vérifiées en particulier
si $\mathfrak X_0$ est une courbe simple sur $k$, et le théorème 4 pourra
s'appliquer en particulier dans la « théorie des modules » pour les courbes de
genre donné\ldots{} Indiquons comment on démontre le théorème 4 : on montre
(cf. proposition 3 plus bas) que \textit{(i)} et \textit{(ii)} impliquent que
tout faisceau cohérent sur $\mathfrak X_0$, localement isomorphe au faisceau
fondamental, s'obtient par réduction à partir d'un faisceau de même nature sur
$\mathfrak X$. Partant alors d'un faisceau « ample » sur $\mathfrak X_0$ (il
en existe d'après \textit{(iii)}) on le remonte en un faisceau inversible sur
$\mathfrak X$, et utilisant le théorème 1 on prouve qu'un multiple de ce
dernier définit une immersion de $\mathfrak X$ dans le complété formel d'un
schéma $\mathbf P_A^r$ (« projectif type » de dimension $r$ au-dessus de
$A$).

Pour la démonstration des théorèmes 1 à 3, on se reportera à [1].

\src{199}
\fsection{4}{Application au théorème de connexion et au « Main theorem » de
Zariski}

Soit $f:X\longrightarrow Y$ un morphisme propre de schémas, alors d'après le
théorème de finitude $f_*(\Osh_X)=\mathcal A$ est un faisceau cohérent sur
$Y$, c'est d'ailleurs un faisceau d'algèbres commutatives, donc il correspond
à un $Y$-schéma $g:Y'\longrightarrow Y$ fini au-dessus de $Y$ (défini par la
condition d'être affine au-dessus de $Y$, i.e. l'image inverse d'un ouvert
affine est affine, et que $g_*(\Osh_{Y'})=\mathcal A$). Il est immédiat que
$f$ se factorise alors canoniquement en $f=gf'$, où
$f':X\longrightarrow Y'$ est un morphisme de $X$ dans $Y'$ qui cette fois-ci
est tel que $f'_*(\Osh_X)=\Osh_{Y'}$. Cette factorisation de $f$ est appelée
la \emph{factorisation de Stein} de $f$. Appliquant le premier théorème de
comparaison, corollaire 1, à $f'$ et à la partie de $Y'$ réduite à un point
$y'$, on trouve que $f'^{-1}(y')=X'$ est connexe (autrement, les sections
formelles de $X$ le long de $X'$ ne formeraient pas un anneau local, or le
complété de $f'_*(\Osh_X)_{y'}=\Osh_{Y',y'}$ est local !). On a prouvé :

\result{THÉORÈME}{5 (« Théorème de connexion » de Zariski)} Soit
$f:X\longrightarrow Y$ un morphisme propre, alors $f$ se factorise de façon
unique (à un isomorphisme près) en $f=gf'$, où
$g:Y'\longrightarrow Y$ est fini et $f':X\longrightarrow Y'$ est tel que
$f'_*(\Osh_X)=\Osh_{Y'}$ (d'où
$g_*(\Osh_{Y'})=f_*(\Osh_X)$). Les fibres de $f'$ sont connexes, i.e.
l'ensemble des composantes connexes d'une fibre $f^{-1}(y)$ de $f$ est en
correspondance biunivoque avec l'ensemble des points de $Y'$ au-dessus de $y$,
i.e. l'ensemble des idéaux maximaux dans $f_*(\Osh_X)_y$.

On en déduit immédiatement les variantes habituelles du théorème de connexion.
Notons ici seulement le

\src{200}
\result{COROLLAIRE}{1} Pour qu'un point $x$ de $X$ soit isolé dans sa fibre
$f^{-1}(y)$ il faut et il suffit que la fibre $f'^{-1}(y')$ (où
$y'=f'(x)$) soit réduite à $x$, ou encore que $f'$ induise un isomorphisme
d'un voisinage de $x$ sur un voisinage de $y'$. L'ensemble de ces points est
un ouvert $U$, et $f'$ induit un isomorphisme de $U$ sur un ouvert de $Y'$.

Pour montrer que $f'$ est un isomorphisme local en $x$, on note que $f'$ induit
un isomorphisme $\Osh_{y'}\longrightarrow\Osh_x$, comme on voit grâce à
$f'_*(\Osh_X)=\Osh_{Y'}$, et au fait que ($f'$ étant une application fermée
dont la fibre en $y'$ est réduite à $x$) les $f'^{-1}(V)$ parcourent un
système fondamental de voisinages de $x$ quand $V$ parcourt un système
fondamental de voisinages de $y'$. On en déduit aussitôt le résultat suivant,
dû à Chevalley dans le cas « géométrique ».

\result{COROLLAIRE}{2} Pour que $f$ soit un morphisme fini, il faut et il
suffit qu'il soit propre et que ses fibres soient finies.

S'il en est ainsi, $f'$ est en effet un isomorphisme d'après ce qui précède.

Soit $f:X\longrightarrow Y$ un morphisme non nécessairement propre, mais
supposons que $X$ soit contenu comme ouvert dans un $Y$-schéma propre
$\widetilde f:\widetilde X\longrightarrow Y$ (ce qui sera le cas en
particulier si $f$ est quasi-projectif). Appliquant le corollaire 1, on voit que
$\widetilde f'$ induit un isomorphisme de l'ensemble $U$ des points de $X$
isolés dans leur fibre sur un ouvert de $Y'$ (et que $U$ est bien ouvert). On
en déduit la variante globale suivante du « Main Theorem » de Zariski.

\result{THÉORÈME}{6 (« Main Theorem » de Zariski)} Soit
$f:X\longrightarrow Y$ un morphisme de type fini. Alors l'ensemble $U$ des
points de $X$ qui sont isolés dans leur fibre est ouvert, et si $f$ est
séparé,\footnote{L'erratum de 1962 remplace ici l'hypothèse imprimée « $f$
quasi-projectif » par l'hypothèse plus faible « $f$ séparé », en vertu de SGA
VIII, 6.2 : tout morphisme quasi-fini et séparé est quasi-projectif.} $U$ est
$Y$-isomorphe à une partie ouverte d'un schéma $Y'$ fini au-dessus de $Y$.

Comme un morphisme de type fini est localement affine, et
\textit{a fortiori} localement quasi-projectif, on déduit aussitôt du théorème
6 les variantes usuelles locales du Main Theorem.

\fsection{5}{Application à l'étude cohomologique des morphismes propres et
plats}

Soit $f:X\longrightarrow Y$ un morphisme propre, soit $F$ un faisceau cohérent
sur $X$, $F$ étant supposé $Y$-plat, i.e. les $F_x$ sont des modules plats
sur les anneaux $\Osh_y$ ($y=f(x)$). Cela signifie aussi que pour tout
$y\in Y$, si nous filtrons $F$ le long de la fibre $f^{-1}(y)$ par les
$\mathfrak m_y^nF$ ($\mathfrak m_y$ étant l'idéal maximal de $\Osh_y$) le
gradué associé est isomorphe à
$(F/\mathfrak m_yF)\otimes_{\kappa(y)}\operatorname{gr}(\Osh_y)$, en d'autres
termes on a
\[
  \mathfrak m_y^nF/\mathfrak m_y^{n+1}F
  =F_y\otimes_{\kappa(y)}
  (\mathfrak m_y^n/\mathfrak m_y^{n+1})
\]
pour tout entier $n$, où $X_y$ désigne la fibre $f^{-1}(y)$ considérée comme
schéma propre au-dessus du corps résiduel $\kappa(y)$ de $y$, et $F_y$ le
faisceau $F/\mathfrak m_yF$ induit par $F$ sur $X_y$. Tenant compte de cet
isomorphisme et du théorème 2, on arrive à des majorations et parfois des
computations des $R^qf_*(F)$ au voisinage de $y$, connaissant la cohomologie de
$X_y$ à coefficients dans $F_y$. Le théorème 2 prend ici la forme
\[
  \widehat{R^qf_*(F)_y}
  =\varprojlim_n H^q(X_y,F/\mathfrak m_y^nF).
\]
Nous allons en signaler seulement la conséquence suivante :

\src{201}
\result{PROPOSITION}{1} Soient $f:X\longrightarrow Y$ un morphisme propre,
$F$ un faisceau sur $X$ cohérent et $Y$-plat. Soit $y\in Y$, soit $q$ un
entier, et supposons que $H^q(X_y,F_y)=0$. Alors $R^qf_*(F)$ est nul au
voisinage de $y$, et pour tout $n$, l'homomorphisme naturel
\[
  R^{q-1}f_*(F)_y\longrightarrow
  H^{q-1}(X_y,F/\mathfrak m_y^nF)
\]
est surjectif.

En particulier, si $f$ est un morphisme plat (i.e. $\Osh_X$ est $Y$-plat)
alors tout faisceau cohérent localement libre $F$ sur $X$ est $Y$-plat. Soient
$F,G$ deux tels faisceaux, appliquons la proposition 1 à
$\Homsh_{\Osh_X}(F,G)$ et à $q=1$, on trouve :

\result{THÉORÈME}{7} Soit $f$ un morphisme propre et plat, soient $F,G$ deux
faisceaux cohérents localement libres sur $X$, soit $y\in Y$ et supposons que
\[
  H^1\bigl(X_y,\Homsh_{\Osh_{X_y}}(F_y,G_y)\bigr)=0,
\]
alors tout homomorphisme $u_0:F_y\longrightarrow G_y$ est induit par un
homomorphisme
\[
  u:F|V\longrightarrow G|V,
  \qquad V=f^{-1}(U),
\]
où $U$ est un voisinage de $y$.

\result{COROLLAIRE}{} Si $u_0$ est un isomorphisme (resp. un monomorphisme, un
épimorphisme) alors il en est de même de $u$, pour $U$ assez petit.

En particulier :

\src{202}
\result{COROLLAIRE}{2} Soit $E_0$ un faisceau cohérent localement libre sur
$X_y$ tel que
\[
  H^1\bigl(X_y;\Homsh_{\Osh_{X_y}}(E_0,E_0)\bigr)=0,
\]
alors deux faisceaux localement libres dont des restrictions à $X_y$ sont
isomorphes à $E_0$, sont isomorphes dans un voisinage de $X_y$.

Ainsi :

\result{COROLLAIRE}{3} Supposons $H^1(X_y,\Osh_{X_y})=0$. Alors deux faisceaux
inversibles sur $X$ (i.e. localement isomorphes à $\Osh_X$) dont les
restrictions à $X_y$ sont isomorphes, sont isomorphes dans un voisinage de $X_y$.\footnote{Le texte imprimé ne précise pas le voisinage. On rétablit ici la portée locale du corollaire 2, également utilisée dans la preuve de la proposition 2.}

On en déduit :

\result{PROPOSITION}{2} Soient $Y$ un schéma connexe, $E$ un faisceau cohérent localement libre de rang au moins $2$ sur $Y$, considérons le fibré en espaces projectifs
$X=\mathbf P(E)$ associé à $E$, muni de son faisceau inversible bien connu
$\Osh_X(1)$. Alors tout faisceau inversible $L$ sur $X$ est isomorphe à un
faisceau de la forme
\[
  f^*(L')\otimes\Osh_X(n),
\]
où $L'$ est un faisceau inversible sur $Y$ et $n$ un entier. Ce dernier est
déterminé de façon unique, et $L'$ est déterminé à un isomorphisme près.

Le corollaire 3 précédent prouve que $L$ est isomorphe à un $\Osh_X(n)$ au
voisinage de chaque fibre. Le reste est à peu près formel.

La proposition 2 permet de déterminer les $Y$-morphismes de $X=\mathbf P(E)$
dans un autre fibré projectif. On trouve en particulier :

\result{COROLLAIRE}{} Soit $u$ un automorphisme de $X=\mathbf P(E)$. Alors il
existe un faisceau inversible $L'$ sur $Y$ et un isomorphisme $v$ de $E$ sur
$E\otimes L'$ tel que $u$ soit l'isomorphisme correspondant
\[
  \mathbf P(E)\xrightarrow{\sim}\mathbf P(E\otimes L')=\mathbf P(E);
\]
le couple $(v,L')$ est déterminé à un isomorphisme près.

Soit $\Gamma$ l'ensemble des classes de fibrés inversibles $L'$ sur $Y$ tels
que $E\otimes L'$ soit isomorphe à $E$. Les éléments sont de torsion, car si
$n$ est le rang de $E$, on voit (en prenant les puissances extérieures
$n$-ièmes) que l'on doit avoir $L'^{\otimes n}\simeq\Osh_Y$. Le corollaire peut
alors s'exprimer en disant qu'on a une suite exacte de groupes
\[
  e\longrightarrow
  \operatorname{Aut}(E)/\Gamma(Y,\Osh_Y^*)
  \longrightarrow \operatorname{Aut}_Y(X)
  \longrightarrow \Gamma\longrightarrow e
\]
(d'ailleurs déduite de la suite exacte de cohomologie déduite de la suite
exacte suivante de faisceaux de groupes
\[
  e\longrightarrow \Osh_Y^*
  \longrightarrow \underline{\operatorname{Aut}}(E)
  \longrightarrow \underline{\operatorname{Aut}}_Y(X)
  \longrightarrow e,
\]
où $\Osh_Y^*$ est le faisceau des « unités » de $\Osh_Y$, identifié au centre
de $\underline{\operatorname{Aut}}(E)$).

\src{203}
\fsection{6}{Application à des théorèmes d'existence et d'unicité pour des
faisceaux et schémas sur un anneau $J$-adique complet}

Le théorème 7 donnait un résultat d'unicité sur des faisceaux cohérents
localement libres, en utilisant les théorèmes 1 et 2. Utilisant le théorème 3,
nous déduisons maintenant des théorèmes d'existence de faisceaux, de morphismes
de schémas, ou de schémas. Dans la suite, $A$ désigne un anneau noethérien
local, séparé et complet. La méthode générale consiste toujours à faire des
constructions formelles, ce qui consiste essentiellement à faire de la
géométrie algébrique sur un anneau artinien, et à en tirer des conclusions de
nature « algébrique » en utilisant les trois théorèmes fondamentaux.

\result{PROPOSITION}{3} Soit $\mathfrak X$ un schéma formel propre et plat
au-dessus de $A$, soit $F_0$ un faisceau localement libre sur $X_0$ tel que
\[
  H^2\bigl(X_0,\Homsh_{\Osh_{X_0}}(F_0,F_0)\bigr)=0.
\]
Alors il existe un faisceau localement libre $F$ sur $\mathfrak X$ qui induit
sur $X_0$ un faisceau isomorphe à $F_0$. (Cet $F$ est d'ailleurs unique à un
isomorphisme près si
$H^1(X_0,\Homsh_{\Osh_{X_0}}(F_0,F_0))=0$.)

On construit de proche en proche des faisceaux localement libres $F_n$ sur les
$X_n$, s'induisant l'un l'autre. La construction de $F_n$ rencontre une
obstruction dans
\[
  H^2\bigl(X_0,\Homsh_{\Osh_{X_0}}(F_0,F_0)\bigr)
  \otimes_{A/J}(J^n/J^{n+1}),
\]
qui est donc nulle par hypothèse. Utilisant maintenant le théorème 3, on
trouve :

\result{COROLLAIRE}{1} Soit $X$ un schéma propre et plat au-dessus de $A$, et
soit $F_0$ comme ci-dessus. Alors il existe un faisceau localement libre $F$
sur $X$ induisant sur $X_0$ un faisceau isomorphe à $F_0$. Cet $F$ est
d'ailleurs unique à un isomorphisme près si
$H^1(X_0,\Homsh_{\Osh_{X_0}}(F_0,F_0))=0$.

Soit $X_0$ un schéma de type fini sur le corps $k$, on suppose $X_0$ simple
(par quoi nous entendons absolument simple) sur $k$, mais pas nécessairement
propre sur $k$. Soit $A$ un anneau artinien local de corps résiduel $k$. Nous
nous intéressons à trouver les schémas $X$ plats sur $A$, tels que
$X\otimes_A k=X_0$ (c'est le point de départ de la « théorie des modules » ou
des « variations de structure » de $X_0$). Il revient au même de se donner un
tel $X$, ou sur l'espace topologique $X_0$ un faisceau $\Osh_X$ muni des
structures suivantes :
\begin{enumerate}
  \item[\textit{i.}] C'est un faisceau de $A$-algèbres.
  \item[\textit{ii.}] Il est muni d'un homomorphisme d'augmentation
    $\Osh_X\longrightarrow\Osh_{X_0}$ (compatible avec les structures de
    $A$-algèbres) ; ces données étant assujetties aux conditions suivantes :
    l'augmentation induit un isomorphisme
    $\Osh_X\otimes_A k\xrightarrow{\sim}\Osh_{X_0}$ ; $\Osh_X$ est plat sur
    $A$, i.e. le gradué associé à $\Osh_X$ filtré par les puissances de l'idéal
    maximal $\mathfrak m$ de $A$ est isomorphe à
    $\operatorname{gr}^0(\Osh_X)\otimes_k\operatorname{gr}(A)$, i.e. on a des
    isomorphismes
    \[
      \mathfrak m^n\Osh_X/\mathfrak m^{n+1}\Osh_X
      =\Osh_{X_0}\otimes_k(\mathfrak m^n/\mathfrak m^{n+1}).
    \]
\end{enumerate}
Le fait fondamental est le suivant :

\src{204}
\result{THÉORÈME}{8} Soit $X_0$ un schéma de type fini et simple sur le corps
$k$, supposons $X_0$ affine. Soit $A$ un anneau local artinien de corps
résiduel $k$, alors il existe un $A$-schéma $X$ plat sur $A$ tel que
$X\otimes_A k=X_0$, et deux tels schémas sont nécessairement isomorphes.

Notons que l'isomorphisme en question n'est pas canonique, car $X$ aura en
général des $A$-automorphismes non triviaux induisant l'identité sur $X_0$.
D'autre part, il n'y a pas en général de choix « canonique » d'un $X$
satisfaisant les conditions données, sauf dans le cas où $A$ est une
$k$-algèbre (cas d'égales caractéristiques) où on peut prendre
$X=X_0\otimes_k A$, i.e. $\Osh_X=\Osh_{X_0}\otimes_k A$ (que $X_0$ soit affine
ou non d'ailleurs). Dans le cas d'inégales caractéristiques, j'ignore en
général, quand $X_0$ n'est pas affine, si on peut « remonter » $X_0$ en un $X$
défini sur $A$. Cependant, soit $n$ un entier $>0$, posons
$A_{n-1}=A/\mathfrak m^n$, et supposons qu'on ait remonté $X_0$ en un
$A_{n-1}$-schéma plat $X_{n-1}$ ; on se propose de remonter $X_{n-1}$ en un
$A_n$-schéma plat $X_n$. On sait déjà en vertu du théorème 8 que c'est possible
localement ; d'autre part on vérifie facilement que, si $U_n$ remonte un ouvert
$U_{n-1}$ de $X_{n-1}$, alors le faisceau des groupes d'automorphismes de $U_n$
(induisant l'identité sur $U_{n-1}$) est canoniquement isomorphe à
\[
  \Theta_{X_0/k}\otimes_k(\mathfrak m^n/\mathfrak m^{n+1})
\]
restreint à $U_{n-1}$, et en particulier est commutatif
($\Theta_{X_0/k}$ désignant le faisceau des germes de $k$-dérivations sur
$X_0$). Il s'ensuit facilement que l'on a une obstruction à construire $X_n$
remontant $X_{n-1}$, qui se trouve dans
\[
  H^2(X_0,\Theta_{X_0/k})\otimes_k
  \mathfrak m^n/\mathfrak m^{n+1}.
\]
Par suite :

\result{COROLLAIRE}{1} Soit $X_0$ un schéma de type fini simple sur $k$,
supposons $H^2(X_0,\Theta_{X_0/k})=0$. Alors pour tout anneau local artinien
$A$ de corps résiduel $k$ il existe un $A$-schéma plat $X$ tel que
$X\otimes_A k=X_0$.

D'ailleurs, si on a pu trouver un $X$ plat sur $A$ remontant $X_0$, alors en
vertu du théorème 8, l'ensemble des classes (à un isomorphisme près) de
$A$-schémas plats qui remontent $X_0$ s'identifie à
\[
  H^1\bigl(X_0,\underline{\operatorname{Aut}}(X)\bigr),
\]
où bien entendu $\underline{\operatorname{Aut}}(X)$ désigne le faisceau des
germes d'automorphismes du faisceau de $A$-algèbres $\Osh_X$ compatibles avec
l'augmentation. La filtration de $\Osh_X$ définit une filtration de
$\underline{\operatorname{Aut}}(X)$, le quotient de ce faisceau par le
$n$-ième sous-groupe de la filtration étant
$\underline{\operatorname{Aut}}(X_n)$ ; le gradué associé à cette filtration
est commutatif, et s'identifie à
$\Theta_{X_0/k}\otimes_k\operatorname{gr}(A)$. En particulier, si
$\mathfrak m^{n+1}$ est la première puissance de $\mathfrak m$ qui est nulle,
alors $F^n(\underline{\operatorname{Aut}}(X))$ (le dernier cran de la
filtration) est dans le centre de $\underline{\operatorname{Aut}}(X)$ et est
isomorphe à $\Theta_{X_0/k}\otimes_k\mathfrak m^n$ ; c'est aussi le faisceau
des germes d'automorphismes de $X$ qui induisent l'identité sur
$X_{n-1}=X\otimes_A A/\mathfrak m^n$. Utilisant ces résultats, on obtient
aussitôt les énoncés suivants :

\src{205}
\result{COROLLAIRE}{2} Soit $X_0$ un schéma de type fini simple sur $k$, soit
$A$ un anneau local artinien de corps résiduel $k$ et d'idéal maximal
$\mathfrak m$, on suppose $\mathfrak m^{n+1}=0$, soit
$A_{n-1}=A/\mathfrak m^n$, et enfin soit $X_{n-1}$ un $A_{n-1}$-schéma plat
tel que $X_{n-1}\otimes_A k=X_0$. Alors l'ensemble des classes (à un
isomorphisme près induisant l'identité sur $X_{n-1}$) de $A$-schémas plats
$X_n$ tels que $X_n\otimes_A A_{n-1}=X_{n-1}$ est vide, ou est un espace
principal homogène sous
\[
  H^1(X_0,\Theta_{X_0/k})\otimes_k\mathfrak m^n.
\]

(Noter que, en général, il n'y a pas d'élément origine privilégié dans ce
dernier espace, car il n'y a pas de façon privilégiée de remonter $X_{n-1}$ en
$X_n$).

\result{COROLLAIRE}{3} Soit $X_0$ un schéma de type fini simple sur $k$, et
supposons $H^1(X_0,\Theta_{X_0/k})=0$. Alors pour tout anneau local artinien
$A$ de corps résiduel $k$, il existe au plus (à un isomorphisme près) un
$A$-schéma plat $X$ tel que $X\otimes_A k=X_0$.

Les corollaires 1 et 3 impliquent immédiatement les énoncés, en apparence plus
généraux, obtenus en y supposant seulement que $A$ est un anneau local
noethérien complet de corps résiduel $k$, pourvu qu'on y introduise $X$ comme
un schéma formel sur $A$ :

\result{THÉORÈME}{9} Soient $k$ un corps, $X_0$ un schéma de type fini simple
sur $k$. Pour tout anneau local noethérien complet $A$ de corps résiduel $k$,
soit $F(A)$ l'ensemble des classes (à un isomorphisme près induisant l'identité
sur $X_0$) de schémas formels $\mathfrak X$ sur $A$, de type fini et plats sur
$A$, tels que $\mathfrak X\otimes_A k=X_0$. Avec ces notations :
\begin{enumerate}
  \item[\textit{i.}] Si $H^1(X_0,\Theta_{X_0/k})=0$, alors pour tout $A$,
    $F(A)$ a au plus un élément.
  \item[\textit{ii.}] Si $H^2(X_0,\Theta_{X_0/k})=0$, alors pour tout $A$,
    $F(A)$ a au moins un élément.
\end{enumerate}

\src{206}
\result{COROLLAIRE}{1} Supposons $X_0$ propre sur $k$. Sous la condition
\textit{(i)}, pour tout $A$, il existe au plus (à un isomorphisme près
induisant l'identité sur $X_0$) un schéma $X$ propre et plat sur $A$, tel que
$X\otimes_A k=X_0$.

On utilise le théorème 3, corollaire 2. Par exemple :

\result{COROLLAIRE}{2} Si $X$ est un schéma propre et plat sur $A$ tel que
$X\otimes_A k$ soit isomorphe au schéma projectif type $\mathbf P_k^r$ de
dimension $r$ sur $k$, alors $X$ est isomorphe à $\mathbf P_A^r$.

(On peut aussi déduire ce résultat de la proposition 3, corollaire 1).

\result{COROLLAIRE}{3} Soit $X_0$ un schéma simple et projectif sur $k$,
supposons que l'on ait
\[
  H^2(X_0,\Osh_{X_0})=H^2(X_0,\Theta_{X_0/k})=0.
\]
Alors pour tout $A$, il existe un schéma projectif et plat $X$ sur $A$, tel que
$X\otimes_A k=X_0$.

On conjugue le théorème 9\textit{(ii)} et le théorème 4. En particulier :

\result{COROLLAIRE}{4} Soit $X_0$ le schéma d'une courbe algébrique complète
et simple sur $k$. Alors pour tout anneau local noethérien complet $A$ de corps
résiduel $k$, il existe un « schéma de courbe simple » $X$ sur $A$, tel que
$X\otimes_A k=X_0$.

\result{REMARQUES}{}

1) Les corollaires 3 et 4 sont surtout intéressants si $k$ est de
caractéristique $p\ne0$, en prenant pour $A$ un anneau de valuation discrète de
caractéristique $0$, de corps résiduel $k$ ; par exemple le « plus petit $A$
possible », savoir celui pour lequel $p$ engendre l'idéal maximal. (En fait, à
cause des théorèmes de Cohen, il suffit de connaître les corollaires 3 et 4
pour un tel anneau $A$). Signalons à ce propos que, selon les spécialistes, on
ne sait pas s'il existe des schémas sur un corps $k$ qui ne sont pas réduction
mod $p$ d'un schéma plat défini sur un tel anneau $A$. Du moins les résultats
de ce numéro donnent-ils un moyen d'investigation systématique de cette
question. Il faudrait commencer
par regarder si la première obstruction qu'on a dans
$H^2(X_0,\Theta_{X_0/k})$ est nécessairement nulle.

\emph{Complément de 1962.} --- J.-P. Serre a construit une variété projective
non singulière, de dimension 3, sur un corps algébriquement clos $k$ de
caractéristique $p>0$, qui ne provient pas par réduction d'un schéma propre sur
un anneau local intègre de corps résiduel $k$ et de corps des fractions de
caractéristique nulle. Mumford aurait trouvé un exemple analogue avec une
surface projective non singulière.

2) On notera que le théorème 3, et la technique correspondante, n'est valable
que pour un anneau de base (local pour fixer les idées) complet.

Pour passer de résultats connus pour le complété d'un anneau local à des
résultats correspondants pour cet anneau local lui-même, il faudrait un
quatrième « théorème fondamental », dont l'énoncé définitif reste à trouver.

\src{207}
3) On comparera les résultats de ce numéro (notamment les corollaires 1 et 2
précédents) et du suivant, avec les résultats de Kodaira--Spencer sur la
variation des structures complexes. Utilisant le théorème conjectural auquel
il vient d'être fait allusion, on devrait pouvoir conclure, sous les conditions
du corollaire 1, mais où $A$ ne serait plus supposé complet, qu'il existe un
anneau $A'$ contenant $A$, fini et non ramifié sur $A$, tel que
$X\otimes_A A'$ et $X'\otimes_A A'$ soient $A'$-isomorphes ($X,X'$ étant deux
$A$-schémas propres et plats donnés, tels que
$X\otimes_A k=X'\otimes_A k=X_0$). C'est ce qu'on peut prouver du moins quand
$X_0=\mathbf P_k^r$, en utilisant le corollaire à la proposition 2. En tout
cas, lorsque $H^1(X_0,\Theta_{X_0/k})=0$, on peut prouver que les fibres de
$X,X'$ au-dessus de tout point $y$ de $Y=\operatorname{Spec}(A)$ sont
isomorphes, du moins quand on passe à la clôture algébrique du corps résiduel
$\kappa(y)$. (On a un résultat local, plus fort en apparence, en ne supposant
plus nécessairement $A$ local). En ce qui concerne les « variations de
structure » pour l'espace projectif, signalons encore la question suivante,
suggérée par un problème correspondant de Kodaira--Spencer. Soit $X$ un schéma
propre et plat au-dessus de l'anneau local intègre $A$ de corps des fractions
$K$, de corps résiduel $k$, supposons que $X\otimes_A K$ soit isomorphe à
$\mathbf P_K^r$, est-il vrai que $X\otimes_A k=X_0$ est isomorphe à
$\mathbf P_k^r$ (du moins sur la clôture algébrique de $k$) ? Dans ce problème,
on peut supposer que $A$ est un anneau de valuation discrète complet. On a un
problème analogue quand $X_0$ est une variété abélienne.

\emph{Complément de 1962.} --- Le conférencier est alors moins optimiste sur
les résultats conjecturés ci-dessus. Le problème relatif aux variations de
structure de l'espace projectif a toutefois été résolu affirmativement par
Hironaka, et le problème analogue pour les variétés abéliennes par Koizumi.

\fsection{7}{Application à la « Théorie des modules »}

Le conférencier venant seulement lui-même d'aborder cette théorie, nous serons
obligés de nous limiter à quelques indications. Pour simplifier, nous nous
plaçons sur un corps $k$, i.e. nous travaillons en égale caractéristique, bien
que le théorème 8 permette aussi de traiter le cas général, sans changement
essentiel, semble-t-il. Nous n'avons pas dépassé actuellement le stade
« formel », cependant le conférencier espère pouvoir construire à partir de là
des schémas de modules véritables dans certains cas, et en particulier
construire pour tout entier $g$ un schéma sur les entiers, jouant le rôle d'un
schéma des modules universel pour les courbes simples de genre $g$.

\emph{Complément de 1962.} --- Mumford a construit les schémas des modules
pour les courbes de genre $g$ (cf. SHMI). Le théorème 10 prouve en outre que
les « schémas de Jacobi d'échelon $n$ » de la théorie des Modules sont non
singuliers, et même simples sur $\mathbf Z$.

\src{208}
Reprenons la situation et les notations du théorème 9, en supposant maintenant
que $A$ est une algèbre locale de rang fini sur $k$, supposé algébriquement
clos pour simplifier. Alors $F(A)$ peut être regardé comme un foncteur
covariant de $A$, à valeurs dans la catégorie des ensembles, un homomorphisme
de $k$-algèbres $A\longrightarrow B$ définissant une application
$F(A)\longrightarrow F(B)$, puisque tout $A$-schéma plat $X$ tel que
$X\otimes_A k=X_0$ donne naissance à un $B$-schéma $X\otimes_A B$ ayant les
mêmes propriétés. Supposons qu'on puisse trouver une $k$-algèbre locale
complète noethérienne $\mathcal O$, et un isomorphisme fonctoriel
\begin{equation}
  \tag{*}
  \operatorname{Hom}(\mathcal O,A)\xrightarrow{\sim}F(A)
\end{equation}
(où le premier membre désigne les homomorphismes de $k$-algèbres). On voit
facilement qu'un tel $\mathcal O$ est déterminé à un isomorphisme canonique
près, on appelle alors le spectre formel $\mathfrak Y$ de $\mathcal O$ (i.e.
l'espace topologique réduit à un point muni d'un faisceau d'anneaux
topologiques réduit à $\mathcal O$) le \emph{schéma formel des modules} pour
$X_0$. (Notons qu'il n'existe pas nécessairement). Soit $\mathfrak r$ l'idéal
maximal de $\mathcal O$, pour tout $n$ soit
$\mathcal O_n=\mathcal O/\mathfrak r^{n+1}$ (donc $\mathcal O_0=k$), alors
l'homomorphisme canonique $\mathcal O\longrightarrow\mathcal O_n$ est un
élément de $\operatorname{Hom}(\mathcal O,\mathcal O_n)$, donc définit un
élément de $F(\mathcal O_n)$, i.e. un $\mathcal O_n$-schéma plat $X_n$ dont la
réduction mod $\mathfrak r$ est $X_0$. Ces $X_n$ se déduisent les uns des
autres par extension des scalaires (i.e. ici par réductions), d'où résulte
qu'ils proviennent d'un schéma formel $\mathfrak X$ bien déterminé sur le
schéma formel des modules $\mathfrak Y$ ; $\mathfrak X$ est plat sur
$\mathfrak Y$ et $\mathfrak X_0=X_0$. L'isomorphisme (*) est alors donné,
comme on voit aussitôt, en associant à tout homomorphisme de $k$-algèbres
$\mathcal O\longrightarrow A$ la classe du $A$-schéma
$\mathfrak X\otimes_{\mathcal O}A$ (i.e. à tout morphisme de $k$-schémas
$Y'=\operatorname{Spec}(A)\longrightarrow\mathfrak Y$, on associe le
$Y'$-schéma $\mathfrak X\otimes_{\mathfrak Y}Y'$ obtenu par changement de
base). De plus, on voit que l'isomorphisme (*) et sa description précédente
seront valables encore si on suppose seulement que $A$ est une $k$-algèbre
locale noethérienne et complète (pas nécessairement artinienne). Bien entendu,
comme d'habitude, $\mathcal O$ peut fort bien \textit{a priori} avoir des
éléments nilpotents, et il semble probable qu'il doive exister des cas où
$\mathcal O$ est lui-même artinien, sans être identique à $k$. C'est dire à
quel point le point de vue de Kodaira--Spencer (se bornant à prendre des $A$
qui sont des anneaux réguliers) est inadéquat \textit{a priori} dans le cas
général.

\src{209}
Il reste à donner des conditions suffisantes pour qu'il existe un schéma formel
des modules pour $X_0$, supposé propre sur $k$. De façon générale, il est
facile de donner des conditions nécessaires et suffisantes simples sur un
foncteur $A\longmapsto F(A)$ (de $k$-algèbres locales de rang fini, à valeurs
dans les ensembles) pour qu'il puisse se mettre sous la forme
$\operatorname{Hom}(\mathcal O,A)$ pour $\mathcal O$ convenable. Nous ne les
détaillerons pas ici. Signalons seulement que dans le cas qui nous occupe, ces
conditions imposent des conditions non triviales de nature cohomologique sur
$X_0$, et il semble peu probable qu'elles soient toujours vérifiées, bien que
le conférencier n'ait pas construit de contre-exemple. Il semble plausible par
contre que la condition $H^0(X_0,\Theta_{X_0/k})=0$ soit une condition
suffisante (bien que nullement nécessaire) pour l'existence d'un schéma formel
de modules. Nous nous bornons à énoncer ici un théorème dans un cas
particulièrement simple (dont l'analogue en théorie des espaces analytiques est
bien connu, cf. Kodaira--Spencer), qui s'établit sans difficulté à l'aide des
résultats du numéro précédent :

\result{THÉORÈME}{10} Soit $X_0$ un schéma propre et simple sur le corps $k$,
tel que
\[
  H^0(X_0,\Theta_{X_0/k})=H^2(X_0,\Theta_{X_0/k})=0.
\]
Alors il existe un schéma formel des modules pour $X_0$, correspondant à un
anneau local régulier $\mathcal O$ (i.e. une algèbre de séries formelles sur
$k$).

Comme nous l'avons déjà signalé, il n'est pas vrai en général que le schéma
formel $\mathfrak X$ sur $\mathcal O$ soit algébrisable ; mais on sait que c'est
vrai cependant quand $X_0$ est projectif et $H^2(X_0,\Osh_{X_0})=0$
(théorème 4), par exemple quand $X_0$ est de dimension 1. C'est ce qui donne
quelque support à l'espoir de construire un schéma de modules sur les entiers
pour les courbes de genre donné\ldots{}

Remarquons aussi que des méthodes comme celles exposées dans ce numéro peuvent
s'appliquer dans la construction et l'étude des variétés de Picard, et bien
d'autres constructions encore. Nous y reviendrons prochainement.

\fsection{8}{Application au groupe fondamental}

Les techniques exposées permettent d'aborder l'étude systématique du groupe
fondamental, sur le modèle de la théorie topologique. Les deux premiers
théorèmes énoncés dans ce numéro sont des généralisations de résultats d'un
travail récent de Lang--Serre.

Soit $X$ un schéma, alors un $X$-schéma $X'$ est appelé un revêtement non
ramifié de $X$ si
\begin{enumerate}
  \item[\textit{i.}] $X'$ est fini sur $X$, i.e. est défini par un faisceau
    cohérent d'algèbres $\mathcal A=\mathcal A(X')$ sur $X$ ;
  \item[\textit{ii.}] $\mathcal A$ est un faisceau localement libre sur $X$ ;
  \item[\textit{iii.}] Pour tout $x\in X$, le quotient
    \[
      \mathcal A_x/\mathfrak m_x\mathcal A_x
      =\mathcal A_x\otimes_{\Osh_x}\kappa(x)
    \]
    est une algèbre séparable sur $\kappa(x)$.
\end{enumerate}

\src{210}
Cette notion de revêtement non ramifié (due à Serre et au conférencier) possède
toutes les propriétés élémentaires auxquelles on peut s'attendre
raisonnablement, et dont nous ne dresserons pas la liste. Bornons-nous à dire
qu'elle donne lieu à une théorie de Galois calquée sur la théorie de Galois
classique (et la contenant ; les démonstrations étant plutôt plus simples que
les démonstrations généralement reçues pour cette dernière) et la théorie
galoisienne des revêtements topologiques. De façon précise, appelons
\emph{point géométrique} d'un schéma $X$ un morphisme $a$ du spectre $\xi$
d'un corps algébriquement clos $\Omega$ dans $X$, i.e. la donnée d'une
extension algébriquement close du corps résiduel $\kappa(x)$ d'un point
$x=|a|$ de $X$ (appelé \emph{localité} du point géométrique $a$). Si alors
$X'$ est un revêtement non ramifié de $X$, on peut lui associer l'ensemble
$E_a(X')$ des « points géométriques de $X'$ au-dessus de $a$ », i.e.
l'ensemble des couples formés d'un $x'\in X'$ au-dessus de $x$ et d'un
$\kappa(x)$-homomorphisme $\kappa(x')\longrightarrow\Omega$. On obtient ainsi
(pour $(X,a)$ fixés) un foncteur $F(X,a)$ de la catégorie $R(X)$ des
revêtements non ramifiés $X'$ de $X$ dans la catégorie des ensembles finis. Si
$X$ est connexe, le couple formé par $R(X)$ et $F(X,a)$ possède les propriétés
formelles qu'il faut pour être isomorphe au couple analogue défini par un
groupe topologique compact totalement discontinu $\pi$ (i.e. limite projective
de groupes finis) convenable : on prend la catégorie $C(\pi)$ des ensembles
finis $E$ sur lesquels $\pi$ opère continûment, et le foncteur identique
$F(\pi)(E)=E$ de cette catégorie dans la catégorie des ensembles finis.
D'ailleurs, le groupe $\pi$ est déterminé à un isomorphisme canonique près par
la condition que $(C(\pi),F(\pi))$ soit isomorphe à un couple donné. En
l'occurrence, $\pi$ s'appelle le \emph{groupe fondamental} du schéma connexe
$X$ en le point géométrique $a$, et se désigne par $\pi_1(X,a)$. Si $X$ n'est
pas connexe, on le remplace par la composante connexe de $x=|a|$. Si cependant
$X$ est connexe, alors les groupes $\pi_1(X,a)$ pour deux points géométriques
$a',a''$ de $X$ sont isomorphes (l'isomorphisme étant déterminé à des
automorphismes intérieurs près), de sorte que comme d'habitude on peut choisir
$a$ au mieux de son intérêt, par exemple au point générique de
$X$ supposé irréductible. Bien entendu, $\pi_1(X,a)$ est un foncteur covariant
en le schéma pointé $(X,a)$. Tout énoncé concernant la classification de
revêtements non ramifiés\footnote{Le texte imprimé porte « revêtements
inséparables » ; le paragraphe traite la catégorie $R(X)$ des revêtements non
ramifiés et son équivalence avec les ensembles finis à action continue de
$\pi_1(X,a)$.} peut se traduire alors en langage de théorie des groupes,
suivant le dictionnaire bien connu (à cela près qu'il faut tenir compte du fait
qu'ici on a des groupes topologiques).

\src{211}
Notre but est de développer l'analogue de la suite exacte d'homotopie des
espaces fibrés, relativement à un morphisme propre $f:X\longrightarrow Y$.
Évidemment, faute de savoir ce que sont les groupes d'homotopie supérieurs, on
n'aura que des résultats nécessairement incomplets. Pour pouvoir appliquer les
théorèmes fondamentaux du n\textsuperscript{o} 3, nous devons d'abord expliciter
quelques lemmes élémentaires concernant les schémas sur des corps ou des
anneaux d'Artin (conformément au procédé général !).

\result{LEMME}{1} Soit $(X',a')$ un revêtement non ramifié ponctué associé à
une représentation ponctuée de $\pi_1(X,a)$ dans un ensemble fini $E$ (muni du
point marqué $e$). Alors le morphisme canonique
$\pi_1(X',a')\longrightarrow\pi_1(X,a)$ identifie le premier groupe au
stabilisateur de $e$ dans $\pi_1(X,a)$ (et est par suite injectif).

\result{LEMME}{2} Soit $X$ un schéma algébrique sur le corps $k$, $k'$ une
extension radicielle de $k$, alors tout revêtement non ramifié de
$X\otimes_k k'$ provient par image réciproque (i.e. extension des scalaires)
d'un revêtement non ramifié de $X$, déterminé à un isomorphisme canonique près.

Il résulte en particulier de ces deux lemmes que pour toute extension
algébrique $K$ de $k$, et tout point géométrique $a'$ de
$X'=X\otimes_k K$ se projetant sur le point géométrique $a$ de $X$,
l'homomorphisme fonctoriel
$\pi_1(X',a')\longrightarrow\pi_1(X,a)$ est injectif.

\result{LEMME}{3} Soit $X$ un schéma complet au-dessus d'un anneau local
artinien $A$, tel que $H^0(X,\Osh_X)=A$, soit $X'$ un revêtement non ramifié de
$X$ et soit $A'=H^0(X',\Osh_{X'})$, qui est donc un anneau fini sur $A$
(pouvant \textit{a priori} être ramifié sur $A$). Soient $X_0,X'_0$ les
sous-schémas réduits associés à $X,X'$ (obtenus en divisant par les faisceaux
des éléments nilpotents dans $\Osh_X,\Osh_{X'}$), soit $k$ un sous-corps de
$A/\mathfrak r(A)$ sur lequel $A/\mathfrak r(A)$ soit fini (ainsi $X_0$ est un
schéma algébrique complet sur $k$, et $X'_0$ en est un revêtement non ramifié).
Soit enfin $\Omega$ une extension algébriquement close de $k$, et considérons
le revêtement non ramifié $X'_0\otimes_k\Omega$ de
$X_0\otimes_k\Omega$. Alors les deux conditions suivantes sont équivalentes :
\begin{enumerate}
  \item[\textit{i.}] $X'_0\otimes_k\Omega$ est complètement décomposé sur
    $X_0\otimes_k\Omega$.
  \item[\textit{ii.}] Le morphisme naturel
    $X'\longrightarrow X\otimes_A A'$ est un isomorphisme.
\end{enumerate}
Sous ces conditions, $A'$ est une extension non ramifiée de $A$. Enfin, si
$X'$ est connexe, la condition \textit{(i)} équivaut à la condition suivante
plus faible en apparence :
\begin{enumerate}
  \item[\textit{i bis.}] $X'_0\otimes_k\Omega$ admet une section régulière
    au-dessus de $X_0\otimes_k\Omega$.
\end{enumerate}
Lorsque la condition \textit{(ii)} est vérifiée, nous dirons que le revêtement
non ramifié $X'$ de $X$ est \emph{géométriquement trivial}.

\src{212}
\result{LEMME}{4} Soit $f:X\longrightarrow Y$ un morphisme propre tel que
$f_*(\Osh_X)=\Osh_Y$. Soient $a$ un point géométrique de $X$, $b$ sa
projection sur $Y$. Alors
$\pi_1(X,a)\longrightarrow\pi_1(Y,b)$ est surjectif.

Il faut montrer en effet ceci : si un revêtement non ramifié $Y'$ de $Y$
(correspondant à un faisceau d'algèbres localement libre $\mathcal A$) est tel
que $X\otimes_Y Y'$ est disconnexe, alors $Y'$ l'est aussi. En effet,
$f^*(\mathcal A)$ sera somme directe de deux faisceaux d'anneaux non nuls, donc
son image directe aussi, or cette dernière n'est autre que
$\mathcal A\otimes f_*(\Osh_X)=\mathcal A$.

\result{LEMME}{5} Soit $X$ un schéma complet au-dessus d'un corps $k$, on
suppose que $H^0(X,\Osh_X)$ est un anneau local $A$ et que
$A/\mathfrak r(A)$ est radiciel sur $k$. Soit $\Omega$ une clôture algébrique
de $k$, et soit $\overline X=X\otimes_k\Omega$ (il est connexe). Choisissons un
point géométrique $\overline a$ de $\overline X$, se projetant sur le point
géométrique $a$ de $X$, alors on a une suite exacte
\[
  e\longrightarrow\pi_1(\overline X,\overline a)
  \longrightarrow\pi_1(X,a)
  \longrightarrow\pi_1(k,b)\longrightarrow e
\]
(où $\pi_1(k,b)$ est le groupe de Galois de $\Omega$ sur $k$).

Le fait que le premier homomorphisme est injectif a été vu avec les lemmes 1 et
2, l'exactitude au milieu résulte du lemme 3, enfin la surjectivité du dernier
homomorphisme (qui seul fait appel au fait que $A/\mathfrak r(A)$ soit
radiciel) résulte du lemme 4.

\result{PROPOSITION}{4} Soit $f:X\longrightarrow Y$ un morphisme propre et
plat tel que pour tout $y\in Y$,
$H^0(f^{-1}(y),\Osh_{f^{-1}(y)})$ soit une algèbre séparable sur le corps
résiduel $\kappa(y)$ (ce qui est le cas par exemple si $f^{-1}(y)$ est un
schéma séparable sur $\kappa(y)$, i.e. réduit et tel que les corps correspondant
à ces composantes irréductibles soient des extensions séparables de
$\kappa(y)$). Alors le revêtement
$Y'$ de $Y$ associé à $f_*(\Osh_X)$ est non ramifié.

La démonstration est facile grâce au théorème 2.

\src{213}
Cette proposition, jointe au lemme 1, ramène pratiquement l'étude homotopique
des morphismes propres et plats (à fibres séparables) au cas où on a
$f_*(\Osh_X)=\Osh_Y$ (car utilisant la factorisation de Stein, on remplacera
$Y$ par $Y'$).

\result{REMARQUE}{} Un morphisme plat de type fini dont les fibres sont des
schémas séparables (resp. simples) est appelé \emph{séparable} (resp.
\emph{simple}). On montre que si $f$ est plat et si $f^{-1}(y)$ est séparable
(resp. simple) alors il existe un voisinage de $f^{-1}(y)$ sur lequel $f$ est
séparable (resp. simple). Le même résultat est valable pour « absolument
normal » (« théorème de Bertini »).

Soit $f:X\longrightarrow Y$ un morphisme propre tel que
\begin{equation}
  \tag{i}
  f_*(\Osh_X)=\Osh_Y,
\end{equation}
et soit $X'$ un schéma fini sur $X$. Soit $Y'$ le revêtement de $Y$
correspondant à la factorisation de Stein de $X'\longrightarrow Y$
(cf. théorème 5). Soit $y\in Y$, donc l'ensemble des composantes connexes de la
fibre $F'$ de $X'$ sur $y$ s'identifie à l'ensemble des points $y'$ de $Y'$
au-dessus de $y$ (théorème 5). Considérons le morphisme évident
\begin{equation}
  \tag{*}
  X'\longrightarrow X\times_Y Y'
\end{equation}
déduit des morphismes naturels $X'\longrightarrow X$ et
$X'\longrightarrow Y'$\footnote{Le texte imprimé omet le prime sur le second
$X'$ ; le morphisme vers $Y'$ est celui de la factorisation de Stein de
$X'\to Y$.} ; ce sera un isomorphisme chaque fois que $X'$ est de
la forme $X\times_Y Y''$, où $Y''$ est un revêtement non ramifié de $Y$, et
alors $Y'$ ne sera autre que $Y''$ et (*) sera l'identité. Nous voulons
précisément donner des conditions moyennant lesquelles $X'$ est de la forme
qu'on vient d'indiquer, i.e. que $Y'$ est non ramifié et (*) un isomorphisme.
Pour ceci, introduisons la fibre $F$ de $X$ en $y$, c'est un schéma propre sur
$\kappa(y)$ dont $F'$ est un revêtement (non ramifié si $X'$ l'est). Soit
$F'_1$ une composante connexe de $F'$ correspondant à un point $y'_1$ de $Y'$
au-dessus de $y$. Supposons de plus
\begin{enumerate}
  \item[\textit{(ii)}] $X'$ est non ramifié sur $X$ aux points de $F'_1$ (par
    suite $F'_1$ est un revêtement non ramifié de $F$), et
  \item[\textit{(iii)}] $F'_1$ est un revêtement géométriquement trivial de
    $F$ (cf. lemme 3).
\end{enumerate}

\result{THÉORÈME}{11} Sous ces conditions, il existe un voisinage ouvert $U'$
de $y'_1$ dans $Y'$, tel que (*) soit un isomorphisme au-dessus de $U'$. De
plus $Y'$ est non ramifié en $y'_1$ au-dessus de $Y$ (mais peut être ramifié
aux autres points $y'$ de $Y'$ au-dessus de $y$).

Bien entendu, les conditions \textit{(ii)} et \textit{(iii)} sont aussi
nécessaires pour la validité de la conclusion du théorème. La démonstration du
théorème est facile, à l'aide du lemme 3 et du théorème 2.

\src{214}
\result{COROLLAIRE}{1} Supposons toujours \textit{(i)} satisfait. Pour qu'un
revêtement non ramifié sur $X$ soit isomorphe à l'image réciproque d'un
revêtement non ramifié $Y'$ de $Y$, il faut et il suffit que $X'$ induise sur
chaque fibre $f^{-1}(y)$ un revêtement géométriquement trivial.

Comme d'après le théorème 11, l'ensemble des points de $Y$ pour lesquels cette
condition est remplie est ouvert, il suffit de la vérifier sur les points $y$
qui sont fermés\ldots{} Notons l'énoncé suivant équivalent au corollaire 1 : le
noyau de l'homomorphisme (surjectif d'après le lemme 4)
$\pi_1(X)\longrightarrow\pi_1(Y)$ est le sous-groupe invariant fermé engendré
par les images dans $\pi_1(X)$ des $\pi_1(\overline{f^{-1}(y)})$, où
\[
  \overline{f^{-1}(y)}
  =f^{-1}(y)\otimes_{\kappa(y)}\overline{\kappa(y)},
\]
$\overline{\kappa(y)}$ désignant une clôture algébrique de $\kappa(y)$. On
notera que, faute de pouvoir choisir le même point base pour toutes les fibres,
les homomorphismes $\pi_1(\overline{f^{-1}(y)})\longrightarrow\pi_1(X)$ ne sont
de toutes façons déterminés (une fois choisi un point base pour $X$, et par
suite pour $Y$) que modulo composition par un automorphisme intérieur dans
$\pi_1(X)$.

\result{COROLLAIRE}{2} Sous les conditions générales du théorème 11, supposons
de plus $Y,X,X'$ intègres, et soient $K,L,L'$ leurs corps. Alors il existe une
sous-extension séparable $K'$ de $K$ dans $L'$, linéairement disjointe de $L$,
telle que $L'=LK'$ (d'où $L'=L\otimes_K K'$).

(On applique la dernière partie du lemme 3 à la fibre géométrique de $X$). Le
cas d'application le plus intéressant du théorème 11 est obtenu quand $f$ est
un morphisme séparable. Alors $X'$ est aussi séparable sur $Y$, donc en vertu
de la proposition 4, $Y'$ est non ramifié sur $Y$, donc le deuxième membre
$X\times_Y Y'$ dans (*) est non ramifié sur $X$. On en conclut facilement :

\result{COROLLAIRE}{3} Supposons, en plus de \textit{(i)}, que $f$ soit
séparable. Soit $X'$ un revêtement non ramifié et connexe de $X$ ; pour que
$X'$\footnote{Le texte imprimé porte ici « $X$ » ; le sujet du critère est le
revêtement $X'$ introduit dans la phrase précédente.} soit image inverse d'un
revêtement non ramifié $Y'$ de $Y$, il faut et il suffit que le revêtement
$\overline F'$ induit sur une fibre géométrique
$\overline F=f^{-1}(y)$ admette une section régulière.

On notera qu'il n'était pas nécessaire de supposer que $\overline F'$ soit
géométriquement trivial sur $\overline F$ (ce qui sera vrai \textit{a
posteriori}, bien que \textit{a priori} cette condition soit beaucoup plus
forte). Le corollaire 3 est en fait équivalent à l'énoncé suivant :

\src{215}
\result{COROLLAIRE}{4} Soit $f:X\longrightarrow Y$ un morphisme propre et
séparable tel que $f_*(\Osh_X)=\Osh_Y$, soit $\overline F$ la fibre géométrique
d'un point $y\in Y$, et choisissons un point géométrique dans $\overline F$,
d'où par les morphismes $\overline F\longrightarrow X\longrightarrow Y$ des
points géométriques dans $X,Y$, qui seront pris comme points bases pour les
groupes fondamentaux de $\overline F,X,Y$. Sous ces conditions on a la suite
exacte
\[
  \boxed{
    \pi_1(\overline F)\longrightarrow\pi_1(X)
    \longrightarrow\pi_1(Y)\longrightarrow e
  }.
\]

On en déduit facilement les deux énoncés suivants de Serre--Lang, débarrassés
ici de toute hypothèse de normalité :

\result{COROLLAIRE}{5} Soient $X,Y$ des schémas connexes sur un corps $k$, on
suppose que le schéma réduit $X_{\mathrm{réd}}$ est séparable sur $k$ (ce qui
est automatiquement vrai si $k$ est parfait), et que $X$ ou $Y$ est propre sur
$k$.\footnote{Condition ajoutée dans l'erratum de 1962.} Choisissons un point
géométrique $a$ (resp. $b$) dans $X$ (resp. $Y$) d'où un point géométrique
$c=(a,b)$ dans $X\times_kY$ et un morphisme naturel
\[
  \pi_1(X\times_kY,c)\longrightarrow
  \pi_1(X,a)\times\pi_1(Y,b)
\]
(déduit des morphismes fonctoriels de $\pi_1(X\times Y,c)$ dans $\pi_1(X,a)$
et $\pi_1(Y,b)$). Ce morphisme est injectif, et même bijectif si $k$ est
algébriquement clos.

(La surjectivité dans ce dernier cas étant à peu près triviale). On en déduit,
avec Serre--Lang :

\result{COROLLAIRE}{6} Soit $X$ un schéma algébrique connexe et propre sur un corps
algébriquement clos $k$, soit $K$ une extension algébriquement close de $k$,
alors les groupes fondamentaux de $X$ et $X\otimes_kK$ sont les mêmes, i.e.
tout revêtement non ramifié de ce dernier schéma provient par extension des
scalaires d'un revêtement non ramifié (unique à un isomorphisme près) de $X$.

\result{REMARQUES}{}

1) Utilisant la proposition 4, on voit que l'hypothèse
$f_*(\Osh_X)=\Osh_Y$ dans le corollaire 4 n'est pas essentielle. Dans le cas
général, il faut, au lieu de mettre le groupe unité $e$ après $\pi_1(Y)$,
continuer par
\[
  \pi_0(\overline F)\longrightarrow\pi_0(X)
  \longrightarrow\pi_0(Y)\longrightarrow e.
\]
comme en Topologie algébrique.

\src{216}
2) En général, on ne peut rien dire pour l'instant sur le noyau de
$\pi_1(\overline F)\longrightarrow\pi_1(X)$, qui devrait faire intervenir un
$\pi_2(Y)$. Il semble cependant qu'on doive pouvoir démontrer que
$\pi_1(\overline F)\longrightarrow\pi_1(X)$ est injectif si $Y$ est le spectre
d'un anneau local $A$, en s'appuyant sur le théorème 12 plus bas (qui affirme
qu'il en est ainsi si $A$ est complet).

Le théorème 11 n'utilisait que les théorèmes 1 et 2. Nous allons appliquer
maintenant le théorème 3, en nous appuyant sur le lemme élémentaire suivant :

\result{LEMME}{6} Soit $X$ un schéma, $X_0$ le schéma réduit correspondant
(i.e. où on a tué les éléments nilpotents). Alors tout revêtement non ramifié
$X'_0$ de $X_0$ est induit par un revêtement non ramifié $X'$ de $X$,
déterminé à un isomorphisme canonique près.

Ce lemme, de nature purement locale, joue ici un rôle analogue à celui du
théorème 8, en théorie des modules. Le conjuguant avec le théorème d'existence
(théorème 3), on en déduit ici :

\result{THÉORÈME}{12} Soit $A$ un anneau local noethérien et complet de corps
résiduel $k$. Soit $X$ un schéma propre sur $A$. Alors tout revêtement non
ramifié $X'_0$ de $X_0=X\otimes_A k$ est induit par un revêtement non ramifié
$X'$ de $X$, unique à un isomorphisme près.

En d'autres termes :

\result{COROLLAIRE}{1} Choisissons un point géométrique dans $X_0$ comme point
base pour les groupes fondamentaux de $X_0$ et $X$. Alors l'homomorphisme
canonique $\pi_1(X_0)\longrightarrow\pi_1(X)$ est un isomorphisme.

Appliquons maintenant à $X_0$ le lemme 5 (en supposant que
$H^0(X_0,\Osh_{X_0})=k$ pour simplifier), et remarquons que $A$ étant complet,
ses extensions non ramifiées correspondent aux extensions non ramifiées de son
corps résiduel, i.e. on a $\pi_1(Y)=\pi_1(k)$ (où
$Y=\operatorname{Spec}(A)$). On trouve la suite exacte :
\[
  e\longrightarrow\pi_1(\overline X_0)
  \longrightarrow\pi_1(X)\longrightarrow\pi_1(Y)\longrightarrow e.
\]

\result{COROLLAIRE}{2} Soit $f:X\longrightarrow Y$ un morphisme propre et
plat, soient $y_1$ un point de $Y$ et $y_0$ une spécialisation de $y_1$,
considérons les fibres « géométriques » correspondantes
$\overline X_1$ et $\overline X_0$, on suppose $\overline X_0$ séparable et
connexe (ce qui implique que $\overline X_1$ satisfait aux mêmes conditions).
On peut alors trouver un homomorphisme de groupes
$\pi_1(\overline X_1)\longrightarrow\pi_1(\overline X_0)$, défini à un
automorphisme intérieur près, et cet homomorphisme est surjectif.

\src{217}
On pourrait espérer que cet homomorphisme est toujours bijectif,
malheureusement il n'en est rien en général si $\kappa(y_0)$ est de
caractéristique $>0$. Nous allons obtenir cependant plus bas une majoration du
noyau de cet homomorphisme (du moins si $\overline X_0$ est simple), impliquant
que si $\kappa(y_0)$ est de caractéristique $0$, alors l'homomorphisme
ci-dessus est bijectif (résultat qu'on pourrait aussi prouver par voie
transcendante). Du moins obtenons-nous déjà, en tous cas, une majoration du
$\pi_1$ d'une fibre spéciale à l'aide de celui d'une fibre générique. Utilisant
par exemple le fait qu'une courbe algébrique en caractéristique $p$ se remonte
en une courbe en caractéristique $0$ (théorème 9, corollaire 4), on trouve par
voie transcendante :

\result{COROLLAIRE}{3} Soit $X_0$ le schéma d'une courbe complète simple sur
un corps algébriquement clos de caractéristique quelconque, soit $g$ le genre
de $X_0$, alors $\pi_1(X_0)$ admet $2g$ générateurs topologiques, liés par la
relation bien connue.

On en déduit, par une technique bien connue de sections hyperplanes :

\result{COROLLAIRE}{4} Soit $X$ un schéma projectif simple sur un corps
algébriquement clos de caractéristique quelconque, alors $\pi_1(X)$ admet un
nombre fini de générateurs topologiques.

Cherchons le noyau de l'homomorphisme
$\pi_1(\overline X_1)\longrightarrow\pi_1(\overline X_0)$. Pour ceci, on peut
supposer que $Y$ est le spectre d'un anneau de valuation discrète complet
$V=\Osh_y$ ($y=y_0$). La question est équivalente à la suivante : étant donné
un revêtement non ramifié $\overline X'_1$ de $\overline X_1$ (qu'on peut si on
veut supposer galoisien), déterminer sous quelles conditions il provient d'un
revêtement non ramifié de $X_0$. \textit{A priori}, le revêtement donné
provient, par extension des scalaires, d'un revêtement non ramifié $X'_1$ de
$X_1\otimes_K K'$, où $K'$ est une sous-extension finie de la clôture
algébrique $\overline K$ du corps des fractions $K$ de $V$ ; si
$\overline X'_1$ était galoisienne de groupe $G$, on peut choisir $X'_1$
galoisienne de groupe $G$. On a alors : pour que
$X'_1\otimes_{K'}\overline K=\overline X'_1$ provienne d'un revêtement non
ramifié de $X_0$, il faut et il suffit que l'on puisse trouver une extension
finie $K''$ de $K'$ dans $\overline K$ telle que
$X''_1=X'_1\otimes_{K'}K''$ soit de la forme $X''\otimes_{V''}K''$, où $V''$
est la clôture normale de $V$ dans $K''$, et où $X''$ est un revêtement non
ramifié de $X\otimes_VV''$. Supposons par exemple $X_0$ absolument normale,
d'où résulte que $X\otimes_VV''$ est normale (car plat sur $V''$ et à fibre
spéciale normale), dont le corps des fonctions est identique à celui
$K''(X_1)$ de
\[
  X_1\otimes_KK''
  =(X\otimes_VK)\otimes_KK''
  =X\otimes_VK''
  =(X\otimes_VV'')\otimes_{V''}K''.
\]
Soit $L''=K''(X'_1)$ le corps des fonctions de $X'_1\otimes_{K'}K''$, c'est
une extension finie séparable de $K''(X_1)$, et la condition ci-dessus signifie
aussi que $L''$ est une extension non ramifiée du corps des fonctions de
$X\otimes_VV''$ (i.e. le normalisé de $X\otimes_VV''$ dans $L''$ est non
ramifié sur $X\otimes_VV''$). Il suffit d'ailleurs de vérifier que $L''$ est
non ramifié aux points de la fibre spéciale de $X\otimes_VV''$ (puisqu'il est
non ramifié au-dessus de la fibre générique $X_1\otimes_KK''$). Si maintenant
$X_0$ est simple, il résulte du « théorème de pureté » de Nagata--Zariski qu'il
suffit même de vérifier que $L''$ est non ramifié au-dessus de l'anneau local
$O''$ du point générique de la fibre spéciale de $X\otimes_VV''$, qui est un
anneau de valuation discrète, normalisé dans $K''(X_1)$ de l'anneau local
$O\subset K(X_1)$ du point générique de la fibre spéciale de $X$. Donc on a
obtenu :

\src{218}
\result{COROLLAIRE}{5} Sous les conditions et avec les notations précédentes,
pour que le revêtement non ramifié
$X'_1\otimes_{K'}\overline K$ de
$\overline X_1=X_1\otimes_K\overline K$ provienne d'un revêtement non ramifié
de $\overline X_0$, il faut et il suffit qu'il existe une sous-extension finie
$K''$ de $\overline K/K'$ telle que $K''(X'_1)$ soit non ramifié au-dessus de
l'anneau de valuation discrète $O''\subset K''(X_1)$.

Notons maintenant que $O''$ est normalisé dans $K''(X_1)$ de l'anneau de
valuation discrète $O'\subset K'(X_1)$ (normalisé de $O$ dans $K'(X_1)$), et
que $O'$ contient le normalisé $V'$ de $V$ dans $K'$, une uniformisante $u$ de
$V'$ étant aussi une uniformisante de $O'$. Supposons maintenant $X'_1$
galoisien de groupe de Galois $G$ d'ordre $n$ premier à la caractéristique $p$
de $\kappa(y_0)$ (qui est aussi la caractéristique du corps résiduel de $O'$).
Donc $K'(X'_1)$ est « tamely ramified » au-dessus de $O'$, d'où il résulte
facilement (« lemme de Abhyankar ») que si on lui adjoint une racine $n$-ième
$v$ d'une uniformisante de $O'$, il devient non ramifié au-dessus du normalisé
de $O'$ dans $K'(X_1)(v)$. Or on peut prendre pour $v$ une racine $n$-ième
d'une uniformisante de $V'$, ce qui prouve que la condition du corollaire 5 est
satisfaite. (Cette possibilité d'utiliser le lemme de Abhyankar et le théorème
de pureté m'a été vendue par Serre). Pour exprimer le résultat obtenu,
introduisons pour tout groupe compact totalement discontinu $\pi$ le groupe
quotient $\widetilde\pi$ de $\pi$ par le sous-groupe fermé engendré par ses
$p$-sous-groupes de Sylow, i.e. la limite projective des groupes quotients
discrets de $\pi$ qui sont d'ordre premier à $p$. Avec cette notation, on
obtient :

\result{THÉORÈME}{13} Soit $f:X\longrightarrow Y$ un morphisme propre et plat,
soient $y_1$ un point de $Y$ et $y_0$ une spécialisation de $y_1$, on suppose
$\overline X_0$ connexe et simple. Alors l'homomorphisme
$\widetilde\pi_1(\overline X_1)\longrightarrow
\widetilde\pi_1(\overline X_0)$ déduit de l'homomorphisme surjectif du
théorème 12, corollaire 2, est un isomorphisme.

En d'autres termes :

\src{219}
\result{COROLLAIRE}{1} La classification des revêtements galoisiens non
ramifiés, de groupe de Galois d'ordre premier à la caractéristique $p$ de
$\kappa(y_0)$ est la même pour $\overline X_0$ et pour $\overline X_1$.

En particulier, si $\kappa(y_0)$ est de caractéristique nulle, on obtient par
voie algébrique le fait que
$\pi_1(\overline X_1)\longrightarrow\pi_1(\overline X_0)$ est bijectif.

Signalons enfin que les techniques utilisées donnent aussi le résultat
suivant, plus général que le théorème 13 :

\result{THÉORÈME}{14} Soit $f:X\longrightarrow Y$ un morphisme propre et
simple, soit $D$ un sous-schéma fermé de $X$, simple au-dessus de $Y$, de
codimension 1 en tous ses points. Pour une fibre $Z=f^{-1}(z)$ de $f$, soit
$Z'=Z-Z\cap D$ et soit $\pi_1^t(Z')$ le quotient du groupe fondamental
$\pi_1(Z')$ qui classifie les revêtements non ramifiés de $Z'$ qui sont
« tamely ramified » au-dessus de $\overline{Z\cap D}$. Soient $y_0,y_1$ comme
dans le théorème 13. Alors on a un homomorphisme surjectif (défini à un
automorphisme intérieur près)
\[
  \pi_1^t(\overline X'_1)\longrightarrow\pi_1^t(\overline X'_0),
\]
et l'homomorphisme correspondant
\[
  \widetilde\pi_1^t(\overline X'_1)
  \longrightarrow\widetilde\pi_1^t(\overline X'_0)
\]
est un isomorphisme.

On en déduit des variantes correspondantes des corollaires au théorème 13, et
du corollaire 4 au théorème 12. De même, utilisant le théorème 9, corollaire 3,
on trouve par voie transcendante :

\result{COROLLAIRE}{} Soit $X_0$ le schéma d'une courbe complète simple sur un
corps algébriquement clos de caractéristique quelconque, soit
$S=(s_i)_{1\le i\le n}$ une partie finie à $n$ éléments de $X_0$. Alors
$\pi_1^t(X_0-S)$ admet $2g+n$ générateurs topologiques
$x_i,y_i$ ($1\le i\le g$) et $\sigma_j$ ($1\le j\le n$) liés par la relation
\[
  \left(\prod_i x_i y_i x_i^{-1}y_i^{-1}\right)
  \sigma_1\cdots\sigma_n=1,
\]
où les $\sigma_j$ sont des générateurs des groupes d'inertie correspondants
aux $s_j$. Pour tout groupe fini $G$ d'ordre premier à la caractéristique,
engendré par des éléments $x_i,y_i,\sigma_j$ satisfaisant la relation
précédente, il existe un revêtement galoisien non ramifié de $X_0-S$, de groupe
$G$, ayant des groupes d'inertie aux points $s_j$ engendrés par les
$\sigma_j$.

Lorsque $X_0$ est de genre 0 et que $n=3$, on a une solution du « Problème
des trois points », du moins pour les revêtements galoisiens d'ordre premier à
la caractéristique. (Ici, le théorème 9 est d'ailleurs inutile, et il semble
d'autre part qu'on puisse déduire le corollaire précédent du cas particulier
envisagé dans le problème des trois points).

\src{220}
\result{REMARQUES}{}

1) Une étude plus complète, faisant sans doute intervenir des revêtements
galoisiens généralisés de $X,X_0,X_1$ (à groupe de Galois éventuellement
infinitésimal), devrait permettre de récupérer le noyau dans le théorème 12,
corollaire 2. Par contre, une étude des revêtements admettant de la ramification
pas « tame » semble devoir être beaucoup plus difficile.

2) Le lemme 6, joint à un résultat de Grauert concernant le complété formel
d'un schéma projectif non singulier le long d'une section hyperplane (ou au
théorème, non prouvé pour l'instant, mentionné dans la remarque 2 après le
théorème 11), permet aussi de démontrer en géométrie algébrique « abstraite » le
classique théorème de Lefschetz sur le groupe fondamental.

\begin{center}
  \textbf{BIBLIOGRAPHIE}
\end{center}

\begin{description}
  \item[{[1]}] Dieudonné (J.) et Grothendieck (A.). --- \emph{Éléments de
    géométrie algébrique}, Publications mathématiques de l'Institut des Hautes
    Études Scientifiques (à paraître).
  \item[{[2]}] Grothendieck (Alexander). --- \emph{The cohomology theory of
    abstract algebraic varieties}. International Congress of Mathematicians
    [1958. Edinburgh] (à paraître).
  \item[{[3]}] Serre (Jean-Pierre). --- \emph{Faisceaux algébriques
    cohérents}, Annals of Math., Series 2, t. 61, 1955, p. 197--278.
  \item[{[4]}] Serre (Jean-Pierre). --- \emph{Géométrie algébrique et
    géométrie analytique}, Ann. Institut Fourier Grenoble, t. 6, 1955--1956,
    p. 1--42.
\end{description}

\par\smallskip
\begingroup\footnotesize
\endgroup
